\subsection{Curve pullbacks and Albanese restrictions at the boundary}
\label{subsec:boundary-albanese-restrictions}

The first positive-excess boundary admits further restrictions when the
torsion class comes from a curve, or when the surface has independent
global one-forms.  The first argument retains the cotangent bundle of
the base curve in the conormal calculation.  The second uses ampleness
on the original surface.  Neither argument constructs a fibration from
a curve of square zero.

\begin{proposition}[Curve-pullback excess bound]
\label{prop:boundary-curve-pullback-excess}
Let $k$ be algebraically closed of characteristic $p>0$, let $X$ be a
normal integral projective surface, and let $C\subset X$ be a smooth
integral effective Cartier divisor contained in the smooth locus, with
$d=C^2>0$.  Let $M\ne\OO_X$ have order $p$ and be trivial on $C$.
Suppose its logarithmic form $\omega$ has exact order $pa$ along $C$,
where $a\geq1$ and $p\nmid a$, and put
\[
 \delta=2g(C)-2-pad.
\]
Suppose that $f:X\to B$ is a nonconstant morphism to a smooth integral
projective curve of genus $b\geq1$, and that
$M\simeq f^*P$ for a nontrivial order-$p$ line bundle $P$ on $B$.
Then $C$ dominates $B$. If $n=\deg(f|_C)$, then $p\mid n$, and there is an integer
$0\leq j\leq a-1$ such that
\begin{equation}\label{eq:boundary-curve-pullback-bound}
 \delta\geq\frac{pd}{2(p-1)}+
       \frac{n(b-1)}{(p-1)(j+1)}
 \geq\frac{pd}{2(p-1)}+\frac{n(b-1)}{(p-1)a}.
\end{equation}
In particular, $b\geq2$ excludes the endpoint
$\delta=pd/(2(p-1))$. In characteristic two this reads
$\delta\geq d+n(b-1)/a$, and at $a=1$ it reads
$\delta\geq d+n(b-1)$.  This proposition does not require ampleness of
$C$ or a separate N\'eron--Severi divisibility hypothesis.
\end{proposition}

\begin{proof}
If $f(C)=b_0$ is a point, write the effective Cartier fibre as
$f^*b_0=rC+E$, where $r>0$ and $E$ is effective without $C$ as a
component.  Since $f|_C$ is constant,
\[
 0=(f^*b_0)\cdot C=rd+E\cdot C>0,
\]
a contradiction.  Here intersections with $E$ are nonnegative because
$C$ is contained in the smooth locus.  Thus $f|_C$ is a finite
nonconstant morphism.  Resolve $X$ away from $C$; this preserves $f$,
$M$, and all the curve data.  We may therefore assume $X$ smooth.
The finite map $f|_C$ is flat because both curves are smooth.
Taking the norm of $(f|_C)^*P\simeq\OO_C$ gives $P^n\simeq\OO_B$.
Since $P$ has exact order $p$, it follows that $p\mid n$.

Multiplication by $p$ on the Jacobian of $B$ is surjective on
$k$-points.  Choose $Q\in\operatorname{Pic}^0(B)$ with $Q^p\simeq P$,
and put $L=f^*Q$.  Then $Q$ and $L$ have exact order $p^2$ and
$L^p\simeq M$.  Use a $p$th-power trivialization of $P$ and its pullback
for $M$.  Changing a $p$th-power trivialization on a proper integral
variety only multiplies it by a constant, so this choice does not
change $\omega$ or its order.

Let $r:B'\to B$ be the normalization of the $P$-torsor.  This torsor
is integral: a $p$th root of its defining unit in $k(B)$ would, by
normality, give local unit roots trivializing $P$.  Consequently $r$
is finite and purely inseparable of degree $p$.  The normal projective
curve $B'$ is smooth and is a Frobenius twist of $B$, so $g(B')=b$.
The line bundle $Q'=r^*Q$ has $p$th power trivial.  It is nontrivial, since
otherwise the norm formula for the finite flat morphism $r$ would give
\[
 \OO_B\simeq\operatorname{Nm}_r(r^*Q)\simeq Q^p\simeq P.
\]
Thus its logarithmic form $\theta$ is a nonzero regular one-form on
$B'$, by the injectivity of the logarithmic-form map on a smooth
proper curve recalled in
Lemma~\ref{lem:positive-square-resolved-transfer}.

Let $W\to X$ be the normalization of the $M$-torsor.  Before
normalization this torsor is the base change of the $P$-torsor.
Its map to that torsor factors through $B'$ by normality: pullbacks
of functions integral over the target order are integral, hence
regular, on $W$.  Take the suitable resolution $S\to W$ in
Lemma~\ref{lem:positive-square-resolved-transfer}, and let $D$ be
the smooth strict transform of the inverse curve.  We obtain a
morphism $\varphi:S\to B'$ compatible with $f$ and $r$.
The root-genus and resolution estimates give
\begin{equation}\label{eq:boundary-curve-pullback-square}
 D^2\geq pd-2(p-1)\delta.
\end{equation}
Indeed, the inverse Cartier curve on $W$ has square $pd$ and
normalization genus $g(C)$; its genus defect $\Delta$ satisfies
$2\Delta\leq(p-1)\delta$, while the suitable resolution loses at most
$4\Delta$ in self-intersection.  These are the assertions of
Lemmas~\ref{lem:positive-square-root-genus} and
\ref{lem:positive-square-resolution-loss}.

The map $q:D\to C$ is purely inseparable of degree $p$.  Since $D$
dominates $B'$ and $r\varphi|_D=f|_Cq$, comparison of degrees gives
\[
 \deg(\varphi|_D)=n.
\]
The pullback $L'$ of $L$ to $S$ is $\varphi^*Q'$, and its logarithmic
form is the regular form $\eta=\varphi^*\theta$.

We need the conormal-order calculation without assuming
$D^2>0$ or $\delta=d$.  At the generic point of $C$, the local
root calculation in Lemma~\ref{lem:positive-square-root-genus}
writes
\[
 u=b_0^p+t^{pa}A,\qquad z=b_0+t^ay,\qquad y^p=A,
\]
where $t$ defines $C$, the functions $b_0,z$ are units, and
$d(A|_C)\ne0$.  On the normalized cover $t$ is a uniformizer of
$D$ and $dt,dy$ form a cotangent basis at its generic point.
Write $db_0=R\,dt+S_1\,dA$ on $X$.  Since $p\nmid a$,
\begin{equation}\label{eq:boundary-curve-pullback-differential}
 \eta=\frac{(R+a t^{a-1}y)\,dt+t^a\,dy}{z}.
\end{equation}
The residue $y_0$ belongs to $k(D)=k(C)^{1/p}$ but not to $k(C)$.
Every normalized leading residue of $R$ belongs to $k(C)$.
Thus $R$ and $a t^{a-1}y$ cannot cancel when their orders coincide,
and
\[
 0\leq j:=\operatorname{ord}_D\eta
     =\min\{v_C(R),a-1\}\leq a-1.
\]
The valuation of zero is infinity.  In particular $\eta\ne0$, and
its leading coefficient is purely conormal because $j<a$.

Consider now the entire differential map
\[
 d\varphi:\varphi^*\omega_{B'}\longrightarrow\Omega_S^1.
\]
At the generic point of $D$, the form $\theta$ is a unit multiple
of a generator of $\omega_{B'}$, since $D$ dominates $B'$.
Therefore this map has exact order $j$ along $D$.
Division by $D^j$ gives a regular map
$\varphi^*\omega_{B'}\otimes\OO_S(jD)\to\Omega_S^1$
near the entire curve $D$: divisibility by a power of the smooth
prime divisor is checked at its generic point in the regular local
rings.  Its restriction has zero tangential image generically,
hence everywhere on the integral curve $D$.  The conormal sequence
gives a nonzero map
\begin{equation}\label{eq:boundary-curve-pullback-conormal}
 (\varphi|_D)^*\omega_{B'}\otimes N_{D/S}^{j}
       \longrightarrow N_{D/S}^{-1}.
\end{equation}
Taking degrees yields
\[
 (j+1)D^2+n(2b-2)\leq0.
\]
Together with \eqref{eq:boundary-curve-pullback-square}, this gives
the first inequality in \eqref{eq:boundary-curve-pullback-bound}.
The second follows from $j+1\leq a$ and $b\geq1$.
\end{proof}

For example, in characteristic two with $a=1$ and $b=2$,
the degree $n$ is a positive even integer. The proposition gives
$\delta\geq d+n\geq d+2$. This is a necessary bound and does
not construct a surface realizing it.

\begin{corollary}[Albanese image a curve]
\label{cor:boundary-albanese-curve}
In characteristic two, under the hypotheses of
Proposition~\ref{prop:boundary-curve-pullback-excess} concerning
$k,X,C,M,a$, suppose instead that $X$ is smooth, $M$ is algebraically
trivial, and the Albanese image of $X$ is one-dimensional.  Let
$f:X\to B$ be the Stein factorization of the Albanese map, and put
$b=g(B)$ and $n=\deg(f|_C)$, retaining
$d=C^2$ and $\delta=2g(C)-2-2ad$.  Then $b\geq1$ and
\[
 \delta\geq d+\frac{n(b-1)}a.
\]
If $\delta=d$, then $b=1$ and $\operatorname{Alb}(X)$ is an elliptic
curve with a nonzero two-torsion point over $k$.
\end{corollary}

\begin{proof}
Write the pointed Albanese map as $\alpha=h\circ f:X\to B\to A$.
The normal projective curve $B$ is smooth.  Its map $h$ induces
$u:\operatorname{Jac}(B)\to A$, while the composite
$X\to B\to\operatorname{Jac}(B)$ induces
$v:A\to\operatorname{Jac}(B)$ by the Albanese universal property.
The identities on $X$ show $uv=1_A$.  Surjectivity of $f$, followed
by the Jacobian universal property, gives $vu=1_{\operatorname{Jac}(B)}$.
Hence $A\simeq\operatorname{Jac}(B)$ and $b\geq1$.

Picard--Albanese duality identifies $f^*$ with an isomorphism
\[
 \operatorname{Pic}^0(B)\simeq
       (\operatorname{Pic}^0_{X/k})_{\mathrm{red}}
\]
\cite[Expos\'e~236, Corollary~3.2 and Theorem~3.3(iii),
pp.~236--238]{FGAPicard}.
Every $k$-point of $\operatorname{Pic}^0_{X/k}$ lies in its reduction.
Thus $M=f^*P$ for a unique nontrivial order-two line bundle $P$ on
$B$.  Apply Proposition~\ref{prop:boundary-curve-pullback-excess}.
\end{proof}

\begin{proposition}[Ample wedge restriction]
\label{prop:boundary-ample-wedge}
Let $X$ be a smooth integral projective surface over an algebraically
closed field of characteristic two, let $C\subset X$ be a smooth
integral ample divisor, and put $d=C^2$.  Let $M\ne\OO_X$ have
order two, be trivial on $C$, and have logarithmic form $\omega$ of
exact order $2a$ along $C$, where $a\geq1$ is odd.  Suppose
\[
 2g(C)-2=(2a+1)d,
 \qquad
 \mathcal T=\omega_C\otimes\OO_C(C)^{-(2a+1)}.
\]
If a global regular one-form $\beta$ satisfies $\omega\wedge\beta\ne0$,
then
\begin{equation}\label{eq:boundary-ample-wedge-splitting}
 K_X\sim2aC,\qquad
 \Omega_X^1\simeq\OO_X(2aC)\oplus\OO_X,\qquad
 c_2(X)=0,\qquad\mathcal T\simeq\OO_C.
\end{equation}
Consequently, if $\mathcal T$ is nontrivial, all global regular
one-forms are rationally collinear with $\omega$, their pairwise
wedge products vanish, and the Albanese map has generic differential
rank at most one.
\end{proposition}

\begin{proof}
Adjunction gives $K_X\cdot C=2ad$.  If
$\omega\wedge\beta\ne0$, division by the $2a$th power of the
canonical section of $\OO_X(C)$ gives a nonzero section
\[
 \tau\in H^0(X,\omega_X(-2aC)).
\]
Its effective zero divisor has intersection zero with $C$.
Ampleness forces this divisor to vanish.  A nonzero section of a
line bundle on a smooth surface cannot have an isolated zero
without a divisorial zero, since locally a nonunit nonzero function
has a height-one zero.  Thus $\tau$ is nowhere zero and
$K_X\sim2aC$.

The divided form $\omega/s_C^{2a}$ and the form $\beta$ define
\[
 \OO_X(2aC)\oplus\OO_X\longrightarrow\Omega_X^1
\]
with determinant $\tau$.  The determinant is invertible everywhere,
so this map is an isomorphism and $c_2(X)=0$.  Restricting
$K_X\sim2aC$ to $C$ and applying adjunction gives
$\mathcal T\simeq\OO_C$.

If $\mathcal T$ is nontrivial, no such $\beta$ exists.  The nonzero
form $\omega$ spans all global one-forms over $k(X)$, so all their
pairwise wedges vanish.  Invariant global one-forms generate the
cotangent bundle of an abelian variety.  Their pullbacks to $X$
therefore show that the Albanese differential has generic rank at
most one.
\end{proof}

\begin{corollary}[The smallest algebraically trivial endpoint]
\label{cor:boundary-smallest-albanese-gate}
Let $X$ be a smooth integral projective surface over an algebraically
closed field of characteristic two, and let $C\subset X$ be a smooth
integral ample divisor.  Let $M\ne\OO_X$ have order two and be
trivial on $C$.  Suppose its logarithmic form has exact order two,
\[
 2g(C)-2=3C^2,
 \qquad [M]\in2\operatorname{NS}(X).
\]
Then all global regular one-forms are rationally collinear, and the
Albanese map has generic differential rank at most one.  In particular,
it is not generically separable onto a two-dimensional image.

If $M$ is algebraically trivial, its Albanese geometry must be one
of the following:
\begin{enumerate}
 \item the Albanese image is a curve and $\operatorname{Alb}(X)$
       is elliptic;
 \item the Albanese image is a surface and the Albanese map has
       generic differential rank at most one.
\end{enumerate}
These conclusions apply, in particular, when $C$ is very ample.
\end{corollary}

\begin{proof}
Corollary~\ref{cor:positive-excess-characteristic-two-boundary}
makes $\mathcal T=\omega_C\otimes\OO_C(C)^{-3}$ a line bundle of
exact order two.  Apply Proposition~\ref{prop:boundary-ample-wedge}.
A morphism generically separable onto a surface image has generic
differential rank two, so it is excluded.

If $M$ is algebraically trivial, the N\'eron--Severi divisibility
condition is automatic.  The Albanese variety is nonzero, because
the $k$-points of a zero-dimensional reduced neutral Picard scheme
are trivial.  A curve image is handled by
Corollary~\ref{cor:boundary-albanese-curve}; the only other possible
image dimension is two.
\end{proof}

\begin{remark}[Scope and comparison]
\label{rem:boundary-albanese-scope}
The curve-pullback bound uses an existing morphism $X\to B$.
The curve-pullback bound does not require a fibration arising
from the resolved curve of square zero. Proposition~\ref{prop:boundary-resolved-fibration}
does produce a separate fibration on a suitable covering resolution,
using actual triviality of its normal line and the moving lemma.
The elliptic-base case and the inseparable surface-image case in
Corollary~\ref{cor:boundary-smallest-albanese-gate} remain unresolved
by these arguments.  Existence of either configuration remains open.
The nonclosed one-forms from Proposition~\ref{prop:residual-adjoint}
are compatible with rational collinearity: $d(r\omega)$ can be
nonzero although all pairwise wedges vanish.

For comparison, Cotner--Zavyalov prove a Lefschetz theorem for finite
flat commutative coefficients under a Hodge-equivalence or strongly
ample hypothesis \cite[Theorem~3.3.5, pp.~1861--1862]{CotnerZavyalov}.
Applied to $\mu_2$ and a strongly ample curve on a surface, it gives
injectivity of restriction on order-two line bundles.  Strong
ampleness is an additional hypothesis; their Example~3.3.6 concerns
failure for $\alpha_p$ and does not provide an order-two line-bundle
example.  Our bound instead retains $\deg\omega_{B'}$ in
\eqref{eq:boundary-curve-pullback-conormal}.

The remaining inseparable Albanese case belongs to an established
foliation setting: Conti describes the canonical inseparable
degree-two factorization for a characteristic-two surface whose
Albanese image is a surface and whose Albanese map is inseparable
\cite[Lemma~2.22 and Remark~2.23, p.~13]{ContiPositiveSeveriComparison}.
That description does not itself exclude the endpoint here.
\end{remark}
